\documentclass[aspectratio=169,10pt]{beamer}
\usetheme[numbering=fraction,progressbar=frametitle]{metropolis}
\usepackage{fontspec}
\setsansfont{Latin Modern Sans}
\usepackage{amsmath,amssymb,mathtools,graphicx,booktabs}
\usepackage{polyglossia}
\setdefaultlanguage{french}
\setbeamertemplate{navigation symbols}{}
\setbeamerfont{frametitle}{size=\large}
\setbeamerfont{normal text}{size=\normalsize}
\setlength{\parskip}{.45em}
\hypersetup{pdfauthor={},pdfsubject={Lecture slides},pdfkeywords={},pdftitle={Variables aléatoires et dispersion}}
\providecommand{\N}{\mathbb{N}}
\providecommand{\Z}{\mathbb{Z}}
\providecommand{\D}{\mathbb{D}}
\providecommand{\Q}{\mathbb{Q}}
\providecommand{\R}{\mathbb{R}}
\providecommand{\C}{\mathbb{C}}
\providecommand{\vect}[1]{\overrightarrow{#1}}
\providecommand{\vi}{\vec{\imath}}
\providecommand{\vj}{\vec{\jmath}}
\providecommand{\vk}{\vec{k}}
\providecommand{\norm}[1]{\left\lVert #1\right\rVert}
\providecommand{\vc}[2]{\begin{pmatrix}#1\\#2\end{pmatrix}}
\providecommand{\deter}[2]{\det\left(#1,#2\right)}
\providecommand{\abs}[1]{\left\lvert #1\right\rvert}
\providecommand{\dd}{\,\mathrm{d}}
\providecommand{\e}{\mathrm{e}}
\providecommand{\Card}{\operatorname{Card}}
\providecommand{\E}{\operatorname{E}}
\providecommand{\V}{\operatorname{V}}
\providecommand{\degre}{^{\circ}}
\providecommand{\un}[1]{\,\mathrm{#1}}
\providecommand{\sci}[2]{#1\times10^{#2}}
\providecommand{\pgcd}{\operatorname{PGCD}}
\providecommand{\ppcm}{\operatorname{PPCM}}
\title{Variables aléatoires et dispersion}
\subtitle{Cours et exemples}
\date{}
\author{}
\begin{document}
\begin{frame}[plain]
\vfill
{\usebeamerfont{title}\inserttitle\par}
\medskip
{\usebeamerfont{subtitle}\insertsubtitle\par}
\medskip
{\usebeamercolor[fg]{progress bar}\rule{\textwidth}{.5pt}}
\vfill
\end{frame}
\begin{frame}{Le parcours du chapitre}
\textbf{Objectifs.} Construire une loi ; interpréter des événements ; calculer espérance, variance, écart type ; utiliser la linéarité ; simuler des échantillons et calculer leur moyenne observée.

\textbf{Prérequis.} Probabilités, moyenne pondérée, carrés, arbres et listes.

\textbf{Repère annuel.} 19-28 janvier : 8 h
\end{frame}
\begin{frame}{Variable, événement et loi}
Sur un univers fini \(\Omega\), une variable aléatoire réelle X est une fonction qui associe un nombre réel à chaque issue. L’événement \(\{X=a\}\) regroupe les issues donnant la valeur a. Sa loi associe à chaque valeur distincte \(x_i\) la probabilité \(p_i=P(X=x_i)\), avec \(p_i\geq0\) et \(\sum p_i=1\).
\end{frame}
\begin{frame}{Variable, événement et loi}
\begin{columns}[T,onlytextwidth]
\begin{column}{.57\textwidth}
Un dé équilibré donne un gain X égal à 2 pour 1 ou 2, et à \(-1\) sinon. La loi est \(P(X=2)=\dfrac{1}{3}\), \(P(X=-1)=\dfrac{2}{3}\). \(\{X\leq0\}=\{3;4;5;6\}\) : probabilité \(\dfrac{2}{3}\).
\end{column}
\begin{column}{.4\textwidth}
\centering
\includegraphics[width=\linewidth,height=4.5cm,keepaspectratio]{figures/premiere-08-1.pdf}
\end{column}
\end{columns}
\end{frame}
\begin{frame}{Variable, événement et loi}
Les valeurs ne sont pas toujours équiprobables même si les issues le sont. Regrouper les issues qui donnent le même gain.
\end{frame}
\begin{frame}{Variable, événement et loi}
Deux pièces équilibrées sont lancées. X compte les piles. Écrire sa loi et \(P(X\leq1)\).
\end{frame}
\begin{frame}{Espérance et linéarité}
L’espérance est \(E(X)=\sum_i x_ip_i\). Elle décrit la moyenne théorique à long terme, pas un résultat nécessairement possible sur un essai. Pour des variables X et Y sur le même univers fini et des réels a,b : \(E(aX+b)=aE(X)+b\) et \(E(X+Y)=E(X)+E(Y)\), sans hypothèse d’indépendance.
\end{frame}
\begin{frame}{Espérance et linéarité}
Pour la loi précédente \(P(X=2)=\dfrac{1}{3}\), \(P(X=-1)=\dfrac{2}{3}\), \(E(X)=0\) : le jeu est équitable si X est le gain net. Pour \(Y=3X-2\), \(E(Y)=3\times0-2=-2\).
\end{frame}
\begin{frame}{Espérance et linéarité}
Une espérance n’est pas une garantie individuelle ; un jeu peut être perdant à chaque horizon observé malgré une espérance positive.
\end{frame}
\begin{frame}{Espérance et linéarité}
Un gain brut vaut 10 avec probabilité \(\dfrac{1}{5}\) et 0 sinon. Quelle mise rend le jeu équitable ?
\end{frame}
\begin{frame}{Variance, écart type et König-Huygens}
En notant \(\mu=E(X)\), \(V(X)=E((X-\mu)^2)=\sum_i(x_i-\mu)^2p_i\). L’écart type est \(\sigma(X)=\sqrt{V(X)}\). La formule de König-Huygens donne \(V(X)=E(X^2)-E(X)^2\). Une variance est toujours positive ou nulle ; l’écart type a la même unité que X, la variance une unité au carré.
\end{frame}
\begin{frame}{Variance, écart type et König-Huygens}
Pour X de valeurs 0 et 4 avec probabilités \(\dfrac{3}{4}\) et \(\dfrac{1}{4}\), \(E(X)=1\), \(E(X^2)=4\), \(V(X)=3\) et \(\sigma=\sqrt3\). Calcul direct : \((0-1)^2(\dfrac{3}{4})+(4-1)^2(\dfrac{1}{4})=3\).
\end{frame}
\begin{frame}{Variance, écart type et König-Huygens}
\(E(X^2)\) n’est pas \(E(X)^2\). Une variance calculée négative signale une erreur ou une très petite erreur d’arrondi à traiter explicitement.
\end{frame}
\begin{frame}{Variance, écart type et König-Huygens}
Pour X de valeurs \(-1,1\) équiprobables, calculer espérance, variance et écart type.
\end{frame}
\begin{frame}{Transformer et comparer des gains}
Pour a,b réels, \(V(aX+b)=a^2V(X)\) et \(\sigma(aX+b)=|a|\sigma(X)\). Une translation déplace la moyenne sans changer la dispersion ; une multiplication change l’échelle. Deux jeux de même espérance peuvent avoir des dispersions très différentes.
\end{frame}
\begin{frame}{Transformer et comparer des gains}
Un jeu A rapporte toujours 2. Un jeu B rapporte 0 ou 4 équiprobablement. Les deux ont espérance 2 ; variances 0 et 4. Leurs distributions ne sont donc pas interchangeables malgré la même moyenne.
\end{frame}
\begin{frame}{Transformer et comparer des gains}
La variance d’une somme n’est pas, en général, la somme des variances ; aucune règle de ce type n’est utilisée sans hypothèses.
\end{frame}
\begin{frame}{Transformer et comparer des gains}
\(E(X)=3\), \(V(X)=4\). Donner les caractéristiques de \(Y=-2X+5\).
\end{frame}
\begin{frame}{Échantillons et moyenne observée}
Un échantillon de taille n est ici une liste de n réalisations indépendantes d’une même variable X. Sa moyenne observée m varie d’un échantillon à l’autre ; on la compare à \(\mu=E(X)\).
\end{frame}
\begin{frame}{Échantillons et moyenne observée}
Pour X de valeurs 0 et 2 équiprobables, \(\mu=1\). L’échantillon 2, 0, 2, 2, 0, 0, 2, 2, 0, 2 a pour moyenne observée 1,2 ; un autre échantillon de même taille donnera en général une autre moyenne. Une simulation compare plusieurs moyennes observées à \(\mu\) sans rien démontrer.
\end{frame}
\begin{frame}{Échantillons et moyenne observée}
Ne pas recopier un seul tirage n fois ; ne pas confondre la moyenne observée d’un échantillon avec l’espérance \(\mu\).
\end{frame}
\begin{frame}{Échantillons et moyenne observée}
\begin{columns}[T,onlytextwidth]
\begin{column}{.57\textwidth}
X prend les valeurs 0 et 2 de façon équiprobable. Calculer la moyenne observée de l’échantillon 2, 0, 2, 2, 0, puis son écart à \(E(X)\).
\end{column}
\begin{column}{.4\textwidth}
\centering
\includegraphics[width=\linewidth,height=4.5cm,keepaspectratio]{figures/premiere-08-2.pdf}
\end{column}
\end{columns}
\end{frame}
\end{document}
